\documentclass[11pt]{article}
\usepackage[margin=1in]{geometry}
\usepackage[T1]{fontenc}
\usepackage{lmodern,amsmath,amssymb,amsthm,mathtools,booktabs,array}
\usepackage[colorlinks=true,allcolors=blue]{hyperref}
\usepackage{microtype}
\newtheorem{theorem}{Theorem}[section]
\newtheorem{proposition}[theorem]{Proposition}
\newtheorem{lemma}[theorem]{Lemma}
\newtheorem{corollary}[theorem]{Corollary}
\theoremstyle{definition}
\newtheorem{definition}[theorem]{Definition}
\theoremstyle{remark}
\newtheorem{remark}[theorem]{Remark}
\newcommand{\R}{\mathbb R}
\newcommand{\Q}{\mathbb Q}
\newcommand{\Z}{\mathbb Z}
\newcommand{\eps}{\varepsilon}
\newcommand{\HH}{\mathbb H}
\newcommand{\KK}{\mathcal K}
\newcommand{\A}{\mathcal A}
\newcommand{\Hk}{\mathcal H}
\newcommand{\nuE}{\nu_{\eps}}
\newcommand{\relint}{\operatorname{relint}}
\newcommand{\st}{\operatorname{st}}
\title{Repair notes for \texttt{hyper-lean}\\
\large A global positive content, its uniqueness, reciprocity,\\ and a corrected transfer principle}
\author{Draft for the authors of \texttt{algebraic-probability.tex}}
\date{1 October 2026}
\begin{document}
\maketitle

\begin{abstract}
The probability manuscript leaves one central item open: a single
geometric content on a specified event algebra that is finitely additive,
independent of decomposition and strictly positive. We show that this
follows from classical valuation theory (Groemer's extension theorem and
Hadwiger's characterization) together with one short new argument, a
leading-coefficient positivity lemma in the ordered field $\R(\eps)$.
We also prove a uniqueness statement from cube calibration, an
Ehrhart--Macdonald-type reciprocity for relatively open polytopes, an
elementary self-contained version for the box algebra suitable for
immediate Lean formalization, and we record how to repair the transfer
axiom of the older introductory manuscript. Classical inputs are cited,
not reproved; everything else is proved here and should be checked.
\end{abstract}

\section{Setting and classical inputs}
Throughout, $d\geq1$, $\HH=\R(\eps)$ is ordered with $0<\eps<r$ for all
real $r>0$, and every geometric datum is \emph{standard}: a subset of
$\R^d$ defined from real numbers. Fix a compact convex \emph{board}
$\Omega\subseteq\R^d$ with nonempty interior. Let $\KK(\Omega)$ be the
nonempty compact convex subsets of $\Omega$, and let $\A(\Omega)$ be the
Boolean algebra of subsets of $\Omega$ generated by $\KK(\Omega)$
(complements taken in $\Omega$). Elements of $\A(\Omega)$ are our events.
It contains points, segments, closed disks, closed balls, polytopes, all
their finite unions, and all set differences such as half-open segments,
a disk minus a diameter, or relatively open polytopes.

\paragraph{Sign rule in $\HH$.} A nonzero element
$x=\sum_{i\geq i_0}a_i\eps^i$ (finite sum, $a_{i_0}\neq0$) has the sign of
$a_{i_0}$. All contents below are of this polynomial form.

Write $V_0,\dots,V_d$ for the intrinsic volumes, normalized so that
$V_k(K)$ is the $k$-dimensional volume $\Hk^k(K)$ when $\dim K=k$, and
$V_j(K)=0$ for $j>\dim K$. $V_0$ is the Euler characteristic on convex
sets ($V_0\equiv1$ on $\KK$).

\begin{itemize}
\item[\textbf{(G)}] \emph{Groemer's extension theorem}
\cite[Ch.~6]{schneider}, \cite{klainrota}. Each $V_j$ extends uniquely
to a linear functional $\varphi_j$ on the real vector space $S(\Omega)$
spanned by the indicators $\mathbf 1_K$, $K\in\KK(\Omega)$; that is,
$\varphi_j(\mathbf 1_K)=V_j(K)$.
\item[\textbf{(H)}] \emph{Hadwiger's theorem} \cite{hadwiger,klainrota}.
Every real-valued valuation on $\KK(\R^d)$ that is continuous in the
Hausdorff metric and invariant under rigid motions is a real linear
combination of $V_0,\dots,V_d$.
\item[\textbf{(M)}] \emph{McMullen's Euler-type relation}
\cite{mcmullen}. If $\varphi$ is a translation-invariant valuation on
polytopes, homogeneous of degree $j$, then for every polytope $P$,
$\sum_{F}(-1)^{\dim F}\varphi(F)=(-1)^j\varphi(-P)$, the sum over the
nonempty faces $F$ of $P$ including $P$.
\end{itemize}

\section{The global content}

\begin{lemma}[Indicators]\label{lem:ind}
For every $E\in\A(\Omega)$, $\mathbf 1_E\in S(\Omega)$. More precisely,
if $E$ lies in the algebra generated by $K_1,\dots,K_m\in\KK(\Omega)$,
then $E$ is a disjoint union of nonempty \emph{atoms}
\[
 A_T=\Bigl(\Omega\cap\bigcap_{i\in T}K_i\Bigr)\setminus
     \bigcup_{i\notin T}K_i ,\qquad T\subseteq\{1,\dots,m\},
\]
and each atom has the form $C\setminus F$, with $C$ compact convex and
$F$ a finite union of compact convex sets.
\end{lemma}
\begin{proof}
The atoms partition $\Omega$, and every set in the generated algebra is a
union of atoms. Further,
$\mathbf 1_{A_T}=\prod_{i\in T}\mathbf 1_{K_i}\cdot\mathbf 1_\Omega
 \prod_{i\notin T}(\mathbf 1_\Omega-\mathbf 1_{K_i})$; expanding gives a
signed sum of indicators of finite intersections, which are compact
convex or empty. Finally $C=\Omega\cap\bigcap_{T}K_i$ and
$F=\bigcup_{i\notin T}K_i$.
\end{proof}

\begin{definition}[Global content]
For $E\in\A(\Omega)$ put
\[
 \nuE(E)=\sum_{j=0}^{d}\varphi_j(\mathbf 1_E)\,\eps^{\,d-j}\in\HH .
\]
\end{definition}

For $K\in\KK(\Omega)$ this is $\sum_jV_j(K)\eps^{d-j}=\eps^dg_K(\eps^{-1})$,
the rescaled Wills polynomial used in the manuscript.

\begin{lemma}[Low-dimensional unions]\label{lem:lowdim}
Let $F=K_1\cup\dots\cup K_m$ with $K_i$ compact convex and
$\dim K_i\leq k$. Then $\varphi_j(\mathbf 1_F)=0$ for $j>k$ and
$\varphi_k(\mathbf 1_F)=\Hk^k(F)$.
\end{lemma}
\begin{proof}
Inclusion--exclusion gives
$\mathbf 1_F=\sum_{\emptyset\neq U}(-1)^{|U|+1}\mathbf 1_{K_U}$ with
$K_U=\bigcap_{i\in U}K_i$, convex of dimension $\leq k$. Apply
$\varphi_j$ and the normalization of $V_j$: terms with $j>k$ vanish, and
for $j=k$ we obtain $\sum(-1)^{|U|+1}\Hk^k(K_U)$, where $\Hk^k(K_U)=V_k(K_U)$
holds also when $\dim K_U<k$ (both sides are $0$). Since $\Hk^k$ is a
measure, the same inclusion--exclusion sum equals $\Hk^k(F)$.
\end{proof}

\begin{lemma}[Relatively open pieces are fat]\label{lem:fat}
Let $C$ be compact convex of dimension $k$ and let $U\subseteq C$ be
nonempty and open in $C$. Then $\Hk^k(U)>0$.
\end{lemma}
\begin{proof}
Pick $x\in U$ and $r>0$ with $C\cap B(x,r)\subseteq U$. For small $t>0$,
the homothetic copy $x+t(C-x)$ lies in $C$ (convexity) and in $B(x,r)$.
Hence $\Hk^k(U)\geq t^k\Hk^k(C)>0$.
\end{proof}

\begin{theorem}[Global positive content]\label{thm:global}
The map $\nuE:\A(\Omega)\to\HH$ is well defined and
\begin{enumerate}
\item finitely additive: $\nuE(E\cup E')=\nuE(E)+\nuE(E')$ for disjoint
 $E,E'$; in particular it is independent of any decomposition;
\item equal to $\eps^dg_K(\eps^{-1})$ on compact convex $K$;
\item invariant under rigid motions $g$ with $gE\subseteq\Omega$;
\item strictly positive: $\nuE(E)>0$ for every nonempty $E$;
\item of known leading order: if $k$ is the largest dimension of a convex
 set $C$ occurring in a nonempty atom $C\setminus F$ of $E$, then
 $\nuE(E)=\Hk^k(E)\,\eps^{d-k}+O(\eps^{d-k+1})$ with $\Hk^k(E)>0$.
\end{enumerate}
\end{theorem}
\begin{proof}
(1) The definition uses only $\mathbf 1_E$, and
$\mathbf 1_{E\cup E'}=\mathbf 1_E+\mathbf 1_{E'}$ for disjoint sets;
apply linearity of each $\varphi_j$ (Lemma~\ref{lem:ind} and (G)).
(2) is the definition of $\varphi_j$ on $\KK$. (3) The functional
$\mathbf 1_E\mapsto\varphi_j(\mathbf 1_{gE})$ extends $V_j$, which is
motion invariant; uniqueness in (G) gives invariance.

(4) and (5). By (1) and Lemma~\ref{lem:ind} it suffices to treat one
nonempty atom $A=C\setminus F$ with $\dim C=k$; a sum of positive
elements is positive and the leading orders add as stated. Write
$\nuE(A)=\nuE(C)-\nuE(C\cap F)$. The set $C\cap F$ is a finite union of
compact convex sets of dimension $\leq k$, so Lemma~\ref{lem:lowdim} and
the normalization of $V_j$ on $C$ show that all coefficients of
$\eps^{d-j}$ with $j>k$ vanish, while the coefficient of $\eps^{d-k}$ is
\[
 \Hk^k(C)-\Hk^k(C\cap F)=\Hk^k(C\setminus F)>0
\]
by Lemma~\ref{lem:fat}, since $C\setminus F$ is nonempty and open in $C$.
The remaining terms carry higher powers of $\eps$ with real
coefficients, so the sign rule gives $\nuE(A)>0$.
\end{proof}

\begin{corollary}[Regular finitely additive probability]
$P(E)=\nuE(E)/\nuE(\Omega)$ is a finitely additive probability on
$\A(\Omega)$ with $P(E)>0$ for every nonempty event. Hence
$P(\,\cdot\mid E)$ is defined for every nonempty $E$, and every finite
context of the manuscript (square, cube, disk and ball darts) is the
restriction of this one probability to a finite subalgebra. Moreover
$\st P(E)=\lambda(E)/\lambda(\Omega)$, so classical geometric
probability is the standard part.
\end{corollary}
\begin{proof}
Positivity and additivity are Theorem~\ref{thm:global}; $\nuE(\Omega)$
has leading term $\lambda(\Omega)>0$ at $\eps^0$. For the standard part,
divide the two leading coefficients at order $\eps^0$; these are
$\varphi_d(\mathbf 1_E)=\lambda(E)$ (Lebesgue measure is the additive
extension of $V_d$) and $\lambda(\Omega)$.
\end{proof}

\begin{remark}[Where standardness is used]
Only in Lemma~\ref{lem:fat} and the sign rule: coefficients of
$\nuE(E)$ must be real. The manuscript's example $(\eps/2)\eps-\eps^2<0$
uses an infinitesimal length, i.e.\ a nonstandard coefficient; this is
exactly what the theorem excludes. Countable additivity is not available:
in dimension one, $(0,1]=\bigsqcup_n(\frac1{n+1},\frac1n]$, each half-open
piece has content equal to its length (the endpoint terms cancel), and the
partial sums $1-\frac1{N+1}$ never come within $\eps$ of
$\nu_\eps((0,1])=1$, so they do not converge in the order topology
of~$\HH$.
\end{remark}

\section{Uniqueness from cube calibration}

\begin{theorem}[Uniqueness]\label{thm:unique}
Let $W\subseteq\HH$ be a finite-dimensional real subspace and
$\mu:\KK(\R^d)\to W$ a valuation that is Hausdorff continuous and
invariant under rigid motions. If $\mu([0,a]^d)=(a+\eps)^d$ for $d+1$
distinct values $a>0$, then $\mu=\eps^dg_K(\eps^{-1})$ on $\KK(\R^d)$,
and $\mu$ has exactly one additive extension to $\A(\Omega)$, namely
$\nuE$.
\end{theorem}
\begin{proof}
Apply (H) to each coordinate of $\mu$ with respect to a real basis of
$W$: $\mu=\sum_jc_jV_j$ with $c_j\in W$. Since
$V_j([0,a]^d)=\binom dja^j$, the hypothesis reads
$\sum_j\binom djc_ja^j=\sum_j\binom dj\eps^{d-j}a^j$ at $d+1$ points, so
the two polynomials of degree $\leq d$ in $a$ coincide and
$c_j=\eps^{d-j}$. An additive extension is determined on atoms by
Lemma~\ref{lem:ind}, because $\mathbf 1_{A_T}$ is an integer combination
of indicators of convex sets.
\end{proof}

\begin{remark}
The manuscript's own calibration, $\nuE([0,a])=a+\eps$ in dimension one
plus the product rule on boxes, implies the hypothesis. Continuity and
motion invariance are genuine assumptions; without continuity, Hadwiger's
theorem fails.
\end{remark}

\section{Reciprocity for relatively open polytopes}

\begin{theorem}[Infinitesimal Ehrhart--Macdonald reciprocity]\label{thm:recip}
Let $P\subseteq\Omega$ be a polytope of dimension $k$. Then
\[
 \nuE(\relint P)=\sum_{j=0}^k(-1)^{k-j}V_j(P)\,\eps^{d-j}
 =(-1)^k\,\eps^d\,g_P(-\eps^{-1}).
\]
\end{theorem}
\begin{proof}
The relative interiors of the nonempty faces partition $P$. M\"obius
inversion on the face lattice (an Eulerian poset) gives
$\mathbf 1_{\relint P}=\sum_F(-1)^{k-\dim F}\mathbf 1_F$. Apply
$\varphi_j$ and (M) with $\varphi=V_j$, using $V_j(-P)=V_j(P)$:
$\varphi_j(\mathbf 1_{\relint P})=(-1)^k(-1)^jV_j(P)$.
\end{proof}

Checks: an open segment of length $L$ gives $L\eps^{d-1}-\eps^d$, as in
the manuscript; an open square of side $a$ in the plane gives
$a^2-2a\eps+\eps^2=(a-\eps)^2$; an open cube gives $(a-\eps)^3$. In
general a relatively open box is the product $\prod(a_i-\eps)$, the
mirror image of $\prod(a_i+\eps)$ for the closed box.

\section{The box algebra: an elementary version for Lean}
The following is self-contained, needs no valuation theory, and is the
recommended first formalization target.

Let $\mathcal B_d$ be the algebra of subsets of $[0,1]^d$ generated by
closed axis-parallel boxes with real corners. In dimension one define,
for $E\in\mathcal B_1$,
\[
 \nu_1(E)=\lambda(E)+\eps\,\chi(E),\qquad
 \chi(E)=\sum_{x\in[0,1]}\bigl(\mathbf 1_E(x)-\mathbf 1_E(x^+)\bigr),
\]
where $\mathbf 1_E(x^+)$ is the right limit (taken as $0$ at $x=1$); the
sum is finite because $\mathbf 1_E$ is a step function.

\begin{proposition}[One dimension]\label{prop:1d}
$\nu_1$ is finitely additive on $\mathcal B_1$, $\nu_1([a,b])=b-a+\eps$,
$\nu_1(\{a\})=\eps$, $\nu_1((a,b))=b-a-\eps$, and $\nu_1(E)>0$ for every
nonempty $E$.
\end{proposition}
\begin{proof}
Both $\lambda$ and $\chi$ are linear in $\mathbf 1_E$, which gives
additivity with no choice of decomposition. The three values follow by
evaluating the jump sum. If $\lambda(E)>0$ the content is positive by the
sign rule; otherwise $E$ is a finite nonempty set of points and
$\chi(E)=\#E>0$.
\end{proof}

\begin{proposition}[Boxes by Fubini]\label{prop:box}
Define $\nu_d$ on $\mathcal B_d$ inductively by integrating
$x_1\mapsto\nu_{d-1}(E_{x_1})$, a step function, against $\nu_1$, where
$\int f\,d\nu_1=\int f\,d\lambda+\eps\sum_x(f(x)-f(x^+))$. Then $\nu_d$ is
finitely additive, satisfies $\nu_d(\prod I_i)=\prod\nu_1(I_i)$, is
strictly positive on nonempty sets, and coincides with $\nuE$ on
$\mathcal B_d$.
\end{proposition}
\begin{proof}[Proof sketch]
Additivity is linearity of the integral in the integrand. Products
factor by construction. For positivity, every nonempty $E$ is a disjoint
union of nonempty grid cells $\prod I_i$ with each $I_i$ a point or an
open interval; each cell has content $\prod\nu_1(I_i)>0$. Agreement with
$\nuE$ follows from Theorem~\ref{thm:unique} on boxes, or directly since
both are additive and agree on closed boxes.
\end{proof}

\paragraph{Lean status.} Proposition~\ref{prop:1d} is checked in
\path{Hyper/IntervalAlgebra.lean}, for events presented as finite lists of
pairwise disjoint points and open intervals (every element of
$\mathcal B_1$ has such a presentation; that representation lemma is not
formalized). The file proves \texttt{content\_eq\_of\_set\_eq} (two
presentations of the same set have the same content), \texttt{content\_pos},
and the calibration \texttt{closed\_content}; the proof is exactly the
representation-free formula $\nu_1=\lambda+\eps\chi$ above, with
\texttt{volume} from Mathlib and $\chi$ as a jump sum. The axiom audit
reports only \texttt{propext}, \texttt{Classical.choice}, and
\texttt{Quot.sound}. Next step: Proposition~\ref{prop:box} by induction on
$d$, reusing the products in \path{CubicContent.lean}.

\section{Repairing the transfer axiom of the older manuscript}
The older text asserts a transfer principle for $\R(\eps)$. This is
false: the sentence $\forall x\,(x>0\to\exists y\;y\cdot y=x)$ holds in
$\R$ but fails in $\R(\eps)$ at $x=\eps$. Three repairs exist,
increasing in strength.

\begin{enumerate}
\item \textbf{Real closure (semialgebraic transfer).} Let
$\R(\eps)^{\mathrm{rc}}$ be the real closure, realized inside algebraic
Puiseux series $\sum_{q\in\frac1n\Z}a_q\eps^q$. Real closed fields are
model complete (Tarski, A.~Robinson) \cite{marker}, so
$\R\preccurlyeq\R(\eps)^{\mathrm{rc}}$: every first-order statement in
$+,\cdot,<$ with real parameters transfers. This covers all
semialgebraic calculus (polynomial and algebraic functions, the
intermediate and extreme value theorems for them, and derivatives as
standard parts of difference quotients).
\item \textbf{Hahn series with restricted analytic functions.} The
Hahn field $\R((\eps^{\Q}))$ is real closed and, with restricted
analytic functions given by convergent Taylor expansion, is an
elementary extension of $\R_{\mathrm{an}}$ \cite{dmm94}. Thus
$\sin,\cos,\exp$ on bounded arguments transfer. The repository's
rational-exponent HyperList backend embeds in this field, so this route
fits the existing code (\texttt{Hyper/theory/HyperHahn.lean}).
\item \textbf{Exponentiation and full transfer.} Logarithmic-exponential
series give an elementary extension of $\R_{\mathrm{an},\exp}$
\cite{dmm97}. No o-minimal structure can transfer $\sin$ on infinite
arguments, since $\sin$ on $\R$ is not definable there; for that, and
for transfer of arbitrary internal sets, an ultrapower is required, at
the price of a non-canonical $\eps$.
\end{enumerate}
The honest restatement for the introductory manuscript is therefore
route~1: $\HH^{\mathrm{rc}}$ with semialgebraic transfer, plus route~2
for elementary transcendental functions at finite arguments.

\section*{Status summary}
\begin{center}
\begin{tabular}{p{0.55\linewidth}p{0.36\linewidth}}
\toprule
Statement & Status \\
\midrule
Global additive content on $\A(\Omega)$ (Thm.~\ref{thm:global}(1--3)) &
classical, via (G) \\
Strict positivity and leading order (Thm.~\ref{thm:global}(4--5)) &
proved here; short \\
Uniqueness from cube calibration (Thm.~\ref{thm:unique}) & proved here
from (H) \\
Reciprocity $\nuE(\relint P)=(-1)^k\eps^dg_P(-\eps^{-1})$ & proved here
from (M); likely known in another guise \\
Interval algebra (Prop.~\ref{prop:1d}) & Lean-checked \\
Box algebra in dimension $d$ (Prop.~\ref{prop:box}) & elementary; next
Lean target \\
Transfer for $\R(\eps)$ & false; replaced by routes 1--3 \\
Full $\A(\Omega)$ in Lean & needs intrinsic volumes, not in Mathlib \\
\bottomrule
\end{tabular}
\end{center}

\begin{thebibliography}{9}
\bibitem{schneider} R. Schneider, \emph{Convex Bodies: The
Brunn--Minkowski Theory}, 2nd ed., Cambridge University Press, 2014,
Ch.~6.
\bibitem{klainrota} D. A. Klain and G.-C. Rota, \emph{Introduction to
Geometric Probability}, Cambridge University Press, 1997.
\bibitem{hadwiger} H. Hadwiger, \emph{Vorlesungen \"uber Inhalt,
Oberfl\"ache und Isoperimetrie}, Springer, 1957.
\bibitem{mcmullen} P. McMullen, Valuations and Euler-type relations on
certain classes of convex polytopes, \emph{Proc. London Math. Soc.} (3)
35 (1977), 113--135.
\bibitem{marker} D. Marker, \emph{Model Theory: An Introduction},
Springer, 2002, \S3.3.
\bibitem{dmm94} L. van den Dries, A. Macintyre, D. Marker, The elementary
theory of restricted analytic fields with exponentiation, \emph{Ann. of
Math.} 140 (1994), 183--205.
\bibitem{dmm97} L. van den Dries, A. Macintyre, D. Marker,
Logarithmic-exponential power series, \emph{J. London Math. Soc.} 56
(1997), 417--434.
\end{thebibliography}
\end{document}
